\PassOptionsToPackage{dvipsnames,svgnames,table}{xcolor}
\documentclass[11pt]{article}
\usepackage[a4paper,margin=22mm]{geometry}
\usepackage[no-math]{fontspec}
\setmainfont{DejaVu Serif}
\usepackage{amsmath,amssymb,amsthm,mathtools,graphicx,xcolor,hyperref,xurl,xspace,ifthen,enumerate}
\usepackage[shortlabels]{enumitem}
\setlength{\emergencystretch}{4em}
\AtBeginDocument{\special{pdf:mapline pzdr Dingbats <uzdr.pfb}}
\SetMathAlphabet{\mathbf}{normal}{OT1}{cmr}{bx}{n}
\SetMathAlphabet{\mathbf}{bold}{OT1}{cmr}{bx}{n}
\newtheorem{theo}{Theorem}[section]
\newtheorem{lemm}[theo]{Lemma}
\newtheorem{thm}[theo]{Theorem}
\newtheorem{proposition}[theo]{Proposition}
\newtheorem{corollary}[theo]{Corollary}
\newtheorem{definition}[theo]{Definition}
\newtheorem{theorem}[theo]{Theorem}
\newtheorem*{remas*}{Remarks}
\newtheorem{remas}[theo]{Remarks}
\newtheorem*{exam*}{Example}
\newtheorem{ques}[theo]{Question}
\providecommand{\og}{«\nobreakspace}
\providecommand{\fg}{\nobreakspace»}
\newtheorem{lem}[theo]{Lemma}
\newtheorem{lemma}[theo]{Lemma}
\newtheorem{prop}[theo]{Proposition}
\newtheorem{coro}[theo]{Corollary}
\newtheorem{corol}[theo]{Corollary}
\newtheorem*{conj*}{Conjecture}
\newtheorem{conj}[theo]{Conjecture}
\newtheorem{conjecture}[theo]{Conjecture}
\newtheorem{Proposition}[theo]{Proposition}
\newtheorem{theoreme}[theo]{Theorem}
\theoremstyle{definition}
\newtheorem{defi}[theo]{Definition}
\newtheorem{exam}[theo]{Example}
\newtheorem{exem}[theo]{Example}
\newtheorem{remark}[theo]{Remark}
\newtheorem{example}[theo]{Example}
\newtheorem{exams}[theo]{Examples}
\newtheorem{rema}[theo]{Remark}
\newtheorem{nota}[theo]{Notation}
\newtheorem{notation}[theo]{Notation}
\newtheorem*{notas}{Notations}
\newtheorem{result}[theo]{Result}
\newtheorem{question}[theo]{Question}
\newtheorem{prob}[theo]{Problem}
\newtheorem*{rema*}{Remark}
\newtheorem*{defi*}{Definition}
\newtheorem*{theo*}{Theorem}
\newtheorem*{lemm*}{Lemma}
\newtheorem*{prop*}{Proposition}
\newcommand{\enoncename}{}
\newtheorem{corpusenonce}[theo]{\enoncename}
\newenvironment{enonce}[2][]{\renewcommand{\enoncename}{#2}\begin{corpusenonce}}{\end{corpusenonce}}
\newenvironment{enonce*}[2][]{\par\medskip\noindent\textbf{#2.}\itshape}{\par\medskip}
\newenvironment{proofc}[1][\proofname]{\begin{proof}[#1]}{\end{proof}}
\newcommand{\Mc}{\mathcal{M}}
\newcommand{\pointir}{.\ ---\ }
\newcommand{\equalenv}[2]{\expandafter\let\csname#1\expandafter\endcsname\csname#2\endcsname\expandafter\let\csname end#1\expandafter\endcsname\csname end#2\endcsname}
\providecommand{\diagbox}[3][]{\begin{tikzpicture}[x=1em,y=1em,baseline=(current bounding box.center)]\useasboundingbox(0,0) rectangle(3,1.4);\draw(0,1.4)--(3,0);\node at(.5,.3){#2};\node at(2.5,1.1){#3};\end{tikzpicture}}
\providecommand{\eg}{e.g.\ }
\providecommand{\ie}{i.e.\ }
\providecommand{\cf}{cf.\ }
\providecommand{\longthanks}[1]{\par\smallskip\noindent #1\par\smallskip}
\providecommand{\qbinom}[2]{\genfrac{[}{]}{0pt}{}{#1}{#2}}

\providecommand{\mathof}[1]{\mathrm{#1}}

\numberwithin{equation}{section}


\def\endappendix{}


\usepackage{latexsym,amsmath,amssymb,verbatim,tikz,tikz-cd}
\usepackage{mathtools}
\usepackage{graphicx,xcolor}

\usepackage[english]{babel}
\usepackage{amssymb}
\usepackage{amsmath}
\usepackage{amsthm}
\usepackage{graphicx}
\usepackage{csquotes}
\usepackage{bm}
\usepackage{enumitem}
\usepackage{subcaption}
\usepackage{caption}
\renewcommand{\thesubfigure}{\Alph{subfigure}}
\usepackage{dirtytalk}
\providecommand{\nopunct}{\spacefactor1007\relax}
\newtheorem{problem}[theo]{Problem}

\newcommand{\sym}{\mathrm{Sym}}
\newcommand{\qsym}{\mathrm{QSym}}
\newcommand{\Des}{\operatorname{Des}}
\newcommand{\Asc}{\operatorname{Asc}}
\newcommand{\Syt}{\operatorname{SYT}}
\newcommand{\Short}{\operatorname{Short}}
\newcommand{\Loop}{\operatorname{Loop}}
\newcommand{\Stable}{\operatorname{Stable}}
\newcommand{\len}{\operatorname{len}}
\newcommand{\height}{\operatorname{height}}
\newcommand{\core}{\operatorname{core}}
\newcommand{\res}{\operatorname{res}}
\newcommand{\row}{\operatorname{row}}
\newcommand{\insrt}{\operatorname{insert}}
\newcommand{\UP}{\operatorname{UP}}
\newcommand{\DOWN}{\operatorname{DOWN}}


\newcommand{\ZZ}{\mathbb{Z}}
\newcommand{\NN}{\mathbb{N}}
\newcommand{\RR}{\mathbb{R}}
\newcommand{\QQ}{\mathbb{Q}}
\newcommand{\PP}{\mathbb{P}}

\newcommand{\F}{\mathcal{F}}
\newcommand{\Q}{\mathcal{Q}}
\newcommand{\A}{\mathcal{A}}
\newcommand{\M}{\mathcal{M}}
\newcommand{\I}{\mathcal{I}}
\newcommand{\PM}{\mathcal{PM}}
\newcommand{\Pcal}{\mathcal{P}}


\newcommand{\odds}{\mathrm{odds}}

\makeatletter
\newcommand{\proofstep}[1]{
  \par
  \addvspace{\medskipamount}
  \textbf{#1\@addpunct{.}}\enspace\ignorespaces
}

\newcommand\restr[2]{{
  \left.\kern-\nulldelimiterspace 
  #1 
  \vphantom{\big|} 
  \right|_{#2} 
  }}
\makeatother












\usepackage{cleveref}
\crefname{originaltheorem}{original theorem}{original theorems}
\Crefname{originaltheorem}{Original theorem}{Original theorems}
\makeatletter
\newlabel{thm:criterion schur positive young}{{2.5}{source}{}{sourcepointer-1}{}}
\newlabel{thm:criterion schur positive young@cref}{{[originaltheorem][0][]2.5}{[1][0][]source}{}{}{}}
\makeatother
\title{Uniform counting of sparse supersets forces symmetry of the statistic}
\author{Avichai Marmor}
\date{doi:10.5802/alco.351; CC BY 4.0}
\begin{document}
\maketitle
\noindent\textit{Editorial note.} The counted statement is the author's labelled lemma `lem:sufficent condition symmetry two rows` (source lines 533-537); the proof environment begins on the next line after \textbackslash{}end\{lemma\}, so the pairing is by adjacency.\par\medskip
\noindent\textit{Editorial note.} Three context spans are retained: the Schur-positivity criterion, the definition of respecting a composition together with the symmetry criterion, and the folklore lemma on the complementary statistic; every label used inside the retained text is defined inside the package.\par\medskip
\noindent\textit{Editorial note.} The retained preamble is the editorially prepared standalone preamble retained for this source.\par\medskip
\noindent\textit{Editorial note.} No retained line refers to a label, citation or figure outside the retained spans, so the package declares no external reference. The reference number was inferred by counting the theorem-class environments of that section of the pinned source; it is not confirmed by the retained bibliography or by a label-number dump of the original, so it must be read as a navigation pointer only.\par\medskip
\noindent\textit{Editorial scope.} One original result of the pinned publisher article, counted once: the sufficient condition for symmetry from uniform sparse-superset counts (author label lem:sufficent condition symmetry two rows).\par\medskip
\noindent\textit{Allowed standard prerequisites (editorial).}\begin{itemize}
\item Complementary statistic. The complementary statistic is Dbar(a) := [N-1] minus D(a).. The proof replaces symmetry for D by symmetry for Dbar at source line 539, using the retained folklore lemma.
\item Composition classes and the symmetry criterion. A is symmetric iff |A\_Dbar(alpha)| depends only on the equivalence class of the composition alpha under the relation that ignores the order of parts.. It is the criterion applied at source line 549 after the class sizes have been computed.
\item Sparse sets. A subset of [N-1] is sparse if it contains no two consecutive integers.. Sparsity of both D and J is used to run the two cases at source lines 541-547.
\end{itemize}\medskip
\ifdefined\hypertarget\hypertarget{sourcepointer-1}{}\fi
\noindent\textit{Source reference.} Original theorem 2.5: sources/alco-351-matching-schur-2024/publisher-source.tex, source lines 368--368. Original criterion of Adin and Roichman invoked inside the retained folklore lemma at source line 433 and inside the retained proof. It is a navigation pointer: the counted lemma's own proof is retained line for line. The reference number was inferred by counting the theorem-class environments of that section of the pinned source; it is not confirmed by the retained bibliography or by a label-number dump of the original, so it must be read as a navigation pointer only.\par\medskip
\begin{definition}
\label{def:respect composition}
    Let $\A$ be a finite set with a statistic $D:\A \to 2^{[N-1]}$. The set of elements that \emph{respect} a given composition $\alpha \vDash N$, denoted $\A_D(\alpha)$, consists of the elements $a \in \A$ such that $D(a) \subseteq S_\alpha$, where $S_\alpha$ is the set corresponding to the composition $\alpha$. When $D$ is clear from the context, we may write $\A(\alpha)$ instead.
\end{definition}

\begin{lemma}
\label{lem:my symmetry criterion}
    A set $\A$ is symmetric if and only if $|\A(\alpha)| = |\A(\beta)|$ for all $\alpha \sim \beta \vDash N$.
\end{lemma}

\begin{lemma}
\label{lem:symmetric complement is symmetric}
    Let $\A$ be a symmetric set with respect to a statistic $D:\A \to 2^{[N-1]}$. Then $\A$ is also symmetric with respect to the complementary statistic $\bar{D}:\A \to 2^{[N-1]}$ defined by $a \mapsto [N-1]\setminus D(a)$.
\end{lemma}

\begin{lemma}
\label{lem:sufficent condition symmetry two rows}
    Let $\A$ be a finite set with a sparse statistic $D:\A \to 2^{[N-1]}$, and denote $n = \lfloor\frac{N}{2} \rfloor$. Assume that there are constants $b_k \in \NN$ for $0 \le k \le n$, such that for every sparse set $J \subseteq [N-1]$, $|\A(D \supseteq J)| = b_{|J|}$.
    Then $\A$ is symmetric with respect to $D$.
\end{lemma}

\begin{proof}
    Following Lemma~\ref{lem:symmetric complement is symmetric}, it suffices to prove that $\A$ is symmetric with respect to the complementary statistic $\bar{D}$. As we proceed to prove it by Lemma~\ref{lem:my symmetry criterion}, let us find $|\A_{\bar{D}}(\alpha)|$ for all $\alpha = (\alpha_1, \dots, \alpha_\ell) \vDash N$.
    
    If a composition $\alpha$ has $\max_i (\alpha_i) > 2$, 
    then there exists $j \in [N-2]$ such that $j,j+1 \notin S_\alpha$, where $S_\alpha$ is the set corresponding to $\alpha$. By Definition~\ref{def:respect composition}, for every element $a \in \A_{\bar{D}}(\alpha)$ we have $\bar{D}(a) \subseteq S_\alpha$, so $j, j + 1 \notin \bar{D}(a)$. Consequently, $j,j+1 \in D(a)$, and the set $D(a)$ is not sparse. However,
    $D$ is assumed to be a sparse function, so $\A_{\bar{D}}(\alpha) = \emptyset$.
    
    Now assume that $\max_i (\alpha_i) \le 2$, and denote by $k = |\{i \mid \alpha_i = 2\}|$ the number of occurrences of $2$ in $\alpha$. Thus, we have $\ell = N - k$.
    Consider the set $J = [N-1] \setminus S_\alpha$. Notably, the set $J$ is sparse and consists of $k$ elements.
    Let $a \in \A$ be an element. We have $a \in \A_{\bar{D}}(\alpha)$ if and only if $\bar{D}(a) \subseteq S_\alpha$, or equivalently, $D(a) \supseteq J$. Thus, $|\A_{\bar{D}}(\alpha)| = |\A(D \supseteq J)| = b_k$.
    
    To conclude, if $\alpha$ contains an element larger than $2$ then $\A_{\bar{D}}(\alpha) = \emptyset$. Otherwise, the size of $\A_{\bar{D}}(\alpha)$ depends only on the number of occurrences of $2$ in $\alpha$. Therefore, $|\A_{\bar{D}}(\alpha)| = |\A_{\bar{D}}(\beta)|$ for all $\alpha \sim \beta \vDash N$. By Lemma~\ref{lem:my symmetry criterion}, we obtain that $\A$ is symmetric.
\end{proof}

\end{document}
